\section{Conceptos del Modelo Entidad – Relación}
\subsection*{a. }

\addcontentsline{toc}{subsection}{d.}
    
\textbf{¿Qué es un tipo de relación? Explica las diferencias con respecto a una instancia de relación.}

Primero, una relación es una asociación entre entidades, la cual cuenta con un grado y restricciones de cardinalidad y participación. Un tipo de relación es un subconjunto del producto cartesiano de n-Entidades (por ejemplo, la relación binaria es un subconjunto del producto cartesiano de la entidad A y la entidad B). También la podemos definir como el conjunto de todas las instancias de relación, y precisamente este último punto es la clave sobre cómo diferenciar una instancia de relación y un tipo de relación: una instancia es un elemento en particular del tipo de relación. Por ejemplo, si tenemos la relación binaria Estudiar entre Alumno y Materia, el conjunto de todas las instancias de la relación forman el tipo de relación, mientras que Juanito que estudia Matemáticas es una instancia particular de la relación, donde Juanito forma parte del conjunto de entidades de Alumno y Matemáticas forma parte del conjunto de entidades de Materia.
